\subsection{Multiplicative Networks} \label{subsec:polynomial_ops}

In this section we lay the theoretical foundation to split neural networks that contain non-additive network components as Gating, SE-Blocks or attention modules.
For these components we need to split the stream through matrix multiplication layers and Hadamard product layers.
Moreover, we need to approximate smooth functions (as Softmax and the Sigmoid function) by CPWL functions.

We prove a general Theorem in \cref{app:splitting_multiplicative_maxout_networks}, where we define the framework of so-called rank-$k$ Maxout Multiplication networks.
However, we want to state the Theorem here translated to a more application-focused language restricted to ReLU networks.
Recall that a posynomial is a polynomial with positive coefficients.
Here monotonicity is with respect to the non-negative split input variables. Non-negative multiplication preserves this monotonicity but does not preserve ordinary convexity.

\begin{theorem} \label{thm:main_simplified}
    Any (residual) ReLU network $N$ containing linear, convolutional, maximum, matrix product and Hadamard product components can be split into two rank-$2$ monotone Maxout networks $N^+$ and $N^-$, that each calculate the maximum of a finite set of posynomials in the non-negative split input variables ($N = N^+ - N^-$).
    This construction gives an exact difference representation but does not generally give convex streams. For example, the positive monomial $xy$ is not convex on the positive orthant: it takes value $4$ at the midpoint of $(1,3)$ and $(3,1)$, whereas its value at both endpoints is $3$.
\end{theorem}